% Research Report 3 — How long is the bridge?
% Build: (cd ../lab && python mcrg_analyze.py) && latexmk report-03.tex
\documentclass[11pt,a4paper]{article}

\usepackage{common}
% Generated by lab/mcrg_analyze.py. Do not edit by hand.
\newcommand{\BridgeStepsMax}{6}
\newcommand{\BridgeStepsMin}{4}
\newcommand{\McrgC}{0.5}
\newcommand{\McrgMeasBig}{80}
\newcommand{\McrgMeasSmall}{100}
\newcommand{\McrgNumSmall}{18}


\fancyhf{}
\fancyhead[L]{\small\sffamily Yang--Mills programme \textperiodcentered{} Research Report 3}
\fancyhead[R]{\small\sffamily 2026-09-25}
\fancyfoot[C]{\small\sffamily\thepage}
% Ledger identifiers refer to spec.pdf; they are not links in this document.
\newcommand{\idref}[1]{\textsf{#1}}

\hypersetup{pdftitle={Research Report 3: How long is the bridge?},pdfauthor={Christian Bucher}}

\begin{document}

\begin{center}
{\sffamily\small YANG--MILLS EXISTENCE AND MASS GAP \textperiodcentered{} RESEARCH PROGRAMME}\\[8pt]
{\LARGE\bfseries Research Report 3}\\[4pt]
{\Large How long is the bridge?\\A Monte Carlo renormalization-group measurement}\\[10pt]
{\large Christian Bucher}\footnote{Research and drafting assistance: Claude Code (Anthropic). Numerical results are evidence class E4. The AI tool is not an author~\cite{COPE23}.}\\[4pt]
{\small 2026-09-25 \textperiodcentered{} companion to \texttt{spec.pdf} v1.3 and Research Reports 1--2}
\end{center}

\begin{keybox}[title=Result in brief]
\textbf{This report does not solve the problem and contains no proof.} It measures one number that shapes every proof strategy of Reports~1 and~2: how many factor-2 renormalization-group (RG) steps lead from the scaling regime of the continuum limit into the region where a mass gap is proved.
\begin{itemize}[leftmargin=1.2em]
\item \textbf{Measured:} starting from $\beta=2.2$--$2.6$ ($SU(2)$, $d=4$, Wilson action), the Wilson-projected RG flow reaches the proved region $\beta<1/12$ after \BridgeStepsMin--\BridgeStepsMax{} blocking steps (Table~\ref{tab:flow}).
\item \textbf{Why it matters:} as $a\to0$ the total number of RG steps grows like $\log_2(1/(a\Lambda))$, but all of them except these few take place at weak coupling, where Ba{\l}aban's perturbative RG applies. The non-perturbative core of the problem is therefore a \emph{bounded} number of RG steps at intermediate coupling, independent of $a$. This does not make it easy. It makes it a finite target (\idref{A-047}).
\end{itemize}
\textbf{Status.} Evidence (E4), subject to the projection onto the Wilson action and to finite-volume effects (Section~\ref{sec:limits}). The closure count is unchanged (1 of 9 transitions).
\end{keybox}

% ======================================================================
\section{Why the length of the bridge matters}
Report~1 located the obstruction in a coupling window: rigorous control exists for $\beta<1/12$~\cite{SZZ23}, while the continuum limit lives at $\beta\to\infty$ and its scaling regime begins near $\beta\approx2.2$. Report~2 showed that a proof has to carry the theory across this window by renormalization. A lattice of spacing $a$ needs about $\log_2\bigl(1/(a\Lambda)\bigr)$ factor-2 steps to reach the physical scale. That number diverges as $a\to0$. Most of these steps, however, are at weak coupling, where asymptotic freedom makes the RG perturbatively controllable~\cite{Bal87,Bal88,Bal89a,Bal89b}. The steps that need non-perturbative control are the ones through the window. \emph{If their number is small and independent of $a$, the infinite problem contains a finite non-perturbative core.} This report measures that number.

% ======================================================================
\section{Method}
\textbf{Blocking.} A factor-2 blocking of $SU(2)$ links, of Swendsen type: the blocked link from $2x$ in direction $\mu$ is the projection onto $SU(2)$ of the double link $U_\mu(x)U_\mu(x+\hat\mu)$ plus $c=\McrgC$ times the sum of the six bent paths of length four around it. The blocking is gauge covariant, which the self-test in \texttt{lab/mcrg.py} checks.

\textbf{Matching.} For each coupling $\beta$ we simulate a $16^4$ lattice (\McrgMeasBig{} measurements) and block it once and twice. We also simulate $8^4$ lattices on a grid of \McrgNumSmall{} couplings $\beta'$ (\McrgMeasSmall{} measurements each), blocked once. The RG step $\beta\mapsto\beta'$ is defined by equating an observable of the $n$-times blocked $16^4$ lattice at $\beta$ with the same observable of the $(n-1)$-times blocked $8^4$ lattice at $\beta'$. Both lattices then have the same final size, which cancels finite-volume effects to leading order. We match the plaquette $P$ and the $1\times2$ rectangle $R$ at levels $n=1$ ($8^4$) and $n=2$ ($4^4$). Agreement between the four estimates tests how well the blocked theory is described by a single Wilson coupling.

% ======================================================================
\section{Results}

\begin{table}[h]
\centering\small
\begin{tabular}{@{}llllll@{}}
\toprule
$\beta$ & $\beta'(P,1)$ & $\beta'(R,1)$ & $\beta'(P,2)$ & $\beta'(R,2)$ & $\Delta\beta$\\
\midrule
% Generated by lab/mcrg_analyze.py. Do not edit by hand.
1.0 & $0.397(1)$ & $0.311(6)$ & $0.334(30)$ & $0.326(206)$ & $0.666$\\
1.8 & $0.850(2)$ & $0.685(5)$ & $0.616(46)$ & $0.358(189)$ & $1.184$\\
2.2 & $1.721(2)$ & $1.572(3)$ & $1.605(9)$ & $1.529(36)$ & $0.595$\\
2.6 & --- & $2.480(2)$ & $2.415(3)$ & $2.403(4)$ & $0.185$\\
\bottomrule

\end{tabular}
\caption{Matched couplings after one factor-2 RG step; $(X,n)$ names the observable and the blocking level; errors in the last digits come from the statistical error of the target observable. $\Delta\beta=\beta-\beta'$ uses the most blocked plaquette estimate. For comparison, the two-loop asymptotic value for $SU(2)$ is $\Delta\beta=8b_0\ln2\approx0.26$, and the flow speeds up towards strong coupling. A dash means the target lies outside the simulated grid.}\label{tab:mcrg}
\end{table}

\begin{table}[h]
\centering\small
\begin{tabular}{@{}p{10.2cm}cc@{}}
\toprule
Trajectory $\beta_0\to\beta_1\to\dots$ & steps to $\beta<\tfrac18$ & steps to $\beta<\tfrac1{12}$\\
\midrule
% Generated by lab/mcrg_analyze.py. Do not edit by hand.
2.60 $\to$ 2.42 $\to$ 2.04 $\to$ 1.21 $\to$ 0.41 $\to$ 0.14 $\to$ 0.05 & 6 & 6\\
2.40 $\to$ 2.01 $\to$ 1.14 $\to$ 0.38 $\to$ 0.13 $\to$ 0.04 & 5 & 5\\
2.20 $\to$ 1.60 $\to$ 0.55 $\to$ 0.18 $\to$ 0.06 & 4 & 4\\
\bottomrule

\end{tabular}
\caption{Iterated RG flow. Between measured couplings the step map is interpolated linearly. Below the smallest measured coupling it is continued linearly towards zero, which is conservative, since the strong-coupling flow is faster. $\tfrac18$ is the limit of pointwise curvature methods (Report~1, Proposition~3.2); $\tfrac1{12}$ is the proved threshold~\cite{SZZ23}.}\label{tab:flow}
\end{table}

% ======================================================================
\paragraph{Reading the table.} (1)~At $\beta=2.6$ the once-blocked plaquette lies above every unblocked plaquette in the grid, so the $(P,1)$ matching fails. This shows why matchings at equal blocking depth, $(n,n-1)$ with $n=2$, are the reliable ones: they cancel the blocking artifacts. (2)~At $\beta=2.6$ the step $\Delta\beta\approx0.19$ lies below the two-loop value $0.26$, most likely because the final $4^4$ lattice is small (Section~\ref{sec:limits}). (3)~At $\beta=1.8$, in the middle of the crossover, the four estimates scatter widely. There the blocked theory is poorly described by a single Wilson coupling, which is itself a quantitative sign of the crossover. (4)~Below $\beta=1$ the flow speeds up, as strong coupling predicts.

\section{What this means}

\begin{proposal}[Finite core; recorded as \idref{A-047}, \statusbadge{OPEN}]
Suppose there are a class $\mathcal A$ of lattice gauge actions and an integer $k$ with the following properties. (i)~Ba{\l}aban-type RG control carries the theory at every small spacing $a$ into $\mathcal A$ at a coupling in the scaling regime, uniformly in $a$ and the volume (\idref{A-034}). (ii)~$k$ factor-2 blocking steps map $\mathcal A$ into the region where a mass gap is proved, with gauge-invariant observables under control. Then the continuum mass gap, uniform in physical units, follows from the gap in that region, transported back through finitely many steps.
\end{proposal}

The measurement suggests $k\le\BridgeStepsMax$ for $SU(2)$ in the Wilson projection. Three points keep this honest:
\begin{enumerate}
\item The implication is itself only sketched, as for Proposal~3.1 of Report~1. Transporting a gap back through RG steps requires control of the fluctuation fields at every step.
\item Step~(ii) is the same certification problem as (R3) in Report~1. It is finite and $a$-independent, but it lives at intermediate coupling, where no rigorous method currently works, and it faces the dimensional barrier of certified integration.
\item $k$ is small, but each step acts on an infinite-dimensional class of effective actions, not on a single coupling.
\end{enumerate}

% ======================================================================
\section{Limits of this evidence}\label{sec:limits}
\begin{itemize}
\item \textbf{Projection.} The blocked theory is not exactly a Wilson theory. Matching one coupling is a projection, and differences between the four estimates in Table~\ref{tab:mcrg} measure its quality.
\item \textbf{Blocking choice.} The step map depends on the blocking (here $c=\McrgC$), and so may the number of steps. The parameter was not varied here; doing so is the obvious next robustness check.
\item \textbf{Volume.} The twice-blocked lattices are $4^4$. At the largest $\beta$ the physical box is small, and matching equal final sizes cancels finite-volume effects only to leading order.
\item \textbf{Statistics.} Moderate statistics; the errors quoted are statistical only.
\end{itemize}

% ======================================================================
\appendix
\section{Deutsche Zusammenfassung}
Dieser Bericht löst das Problem nicht. Er misst eine Zahl, die für jeden Beweisweg entscheidend ist: Wie viele Renormierungsschritte (Blockierung um den Faktor 2) führen vom Kontinuumsbereich in den Bereich, in dem ein Mass Gap bewiesen ist? Ergebnis der Monte-Carlo-Renormierungsgruppe für $SU(2)$: \BridgeStepsMin{} bis \BridgeStepsMax{} Schritte. Das ist wichtig, weil alle übrigen Schritte bei schwacher Kopplung stattfinden, wo Ba{\l}abans störungstheoretische Methoden greifen. Der nicht-störungstheoretische Kern des Problems ist damit eine \emph{feste, endliche} Zahl von Schritten bei mittlerer Kopplung, unabhängig vom Gitterabstand. Leicht wird er dadurch nicht, aber er ist ein endliches Ziel. Es bleibt ein numerischer Beleg (E4) mit den genannten Einschränkungen.

\begin{thebibliography}{Bal89b}
\small
\bibitem[SZZ23]{SZZ23} H.~Shen, R.~Zhu, X.~Zhu, A stochastic analysis approach to lattice Yang--Mills at strong coupling, \emph{Commun. Math. Phys.} 400 (2023) 805--851. doi:10.1007/s00220-022-04609-1
\bibitem[Bal87]{Bal87} T.~Ba{\l}aban, Renormalization group approach to lattice gauge field theories I, \emph{Commun. Math. Phys.} 109 (1987) 249--301.
\bibitem[Bal88]{Bal88} T.~Ba{\l}aban, Renormalization group approach to lattice gauge field theories II, \emph{Commun. Math. Phys.} 116 (1988) 1--22. doi:10.1007/BF01239022
\bibitem[Bal89a]{Bal89a} T.~Ba{\l}aban, Large field renormalization I, \emph{Commun. Math. Phys.} 122 (1989) 175--202. doi:10.1007/BF01257412
\bibitem[Bal89b]{Bal89b} T.~Ba{\l}aban, Large field renormalization II, \emph{Commun. Math. Phys.} 122 (1989) 355--392. doi:10.1007/BF01238433
\bibitem[COPE23]{COPE23} Committee on Publication Ethics, \emph{Authorship and AI tools}, COPE position statement (2023).
\end{thebibliography}

\end{document}
